\documentclass[12pt]{article}
\usepackage[utf8]{inputenc}
\usepackage[T1]{fontenc}
\usepackage[margin=1in]{geometry}
\usepackage{setspace}
\usepackage{xcolor}
\usepackage[colorlinks=true, urlcolor=blue!60!black, linkcolor=blue!60!black]{hyperref}
\setlength{\parindent}{0em}
\setlength{\parskip}{0.6em}
\renewcommand{\baselinestretch}{1.0}
\begin{document}
\begin{center}
{\LARGE \textit{Cedar Point Nursery v.\ Hassid}, 594 U.S.\ 139 (2021)\par}
\vspace{0.5em}
{\large Case Brief --- Two-Page Edition\par}
\end{center}

\section{Memory Jogger}
A rule letting union organizers enter farms 120 days a year was a \emph{per se} physical taking of the right to exclude (\href{https://supreme.justia.com/cases/federal/us/594/20-107/}{Justia}).

\section{Detailed Case Facts}
A California regulation gave unions a ``right to take access'' to agricultural employers' property for up to three hours a day during four 30-day periods a year. Cedar Point Nursery, whose trim shed organizers entered at 5:00 a.m.\ with bullhorns, and Fowler Packing challenged it.

\section{Procedural History}
The district court dismissed, applying \textit{Penn Central} because access was not continuous; the Ninth Circuit affirmed, limiting \emph{per se} treatment to round-the-clock access. The Supreme Court reversed.

\section{Judicial Votes}
Reversed 6--3. Roberts wrote for Thomas, Alito, Gorsuch, Kavanaugh, and Barrett; Kavanaugh concurred; Breyer (with Sotomayor and Kagan) dissented.

\section{Holding}
A regulation that appropriates for third parties a right to physically enter private property is a \emph{per se} taking, however intermittent; isolated trespasses, background property-law privileges, and reasonable inspection conditions are excepted.

\section{Analysis of Opinions}
\textbf{Majority:} physical appropriations, unlike use restrictions, require compensation without \textit{Penn Central} balancing; the right to exclude is fundamental; duration affects only compensation; and \textit{PruneYard} involved property open to the public. \textbf{Kavanaugh:} \textit{Babcock \& Wilcox} supports the growers. \textbf{Breyer:} the rule regulates rather than appropriates, and only permanent occupations are \emph{per se} takings.

\section{Examples: Future Applications}
\textbf{Same-Side:} a statute letting nonprofits survey wildlife on fenced timberland. \textbf{Opposite-Side:} food-safety inspections required as a license condition. \textbf{Fence-Sitter:} union access to a supermarket's private sidewalk.

\section{Critique}
The growers sought an injunction although the Clause only requires compensation, leaving the relation to \textit{Knick} unclear; the exceptions are fuzzy; and \textit{PruneYard}'s survival is uncertain.

\section{Key Quotations}
\begin{quote}
``The essential question is not, as the Ninth Circuit seemed to think, whether the government action at issue comes garbed as a regulation (or statute, or ordinance, or miscellaneous decree). It is whether the government has physically taken property for itself or someone else---by whatever means---or has instead restricted a property owner's ability to use his own property.''
\end{quote}
\noindent --- Chief Justice Roberts, for the Court.

\end{document}
